%% kapa: exam=1 round=25 year=2014 period=2 form=A subject=appraisal_law subject_ko=부동산관계법규 seq=29 number=29
%% source: Q-Net 감정평가사 기출문제 (cst00309) artlSeq=5101207 file=extracted/1008/5101207/2014년 제25회 감정평가사 1차시험 기출문제/2교시(A형).hwp
%% conversion: kordoc 4.24.51 -> tools/gen1.py (auto); review=proofread-sonnet; flags=none
\begin{question}{29}
건축법령상 이행강제금에 관한 설명으로 옳지 않은 것은?
\begin{choices}
  \item 이행강제금은 건축허가 대상 건축물뿐만 아니라 건축신고 대상 건축물에 대해서도 부과할 수 있다.
  \item 연면적 85제곱미터 이하의 주거용 건축물인 경우에는 법정 이행강제금의 2분의 1의 범위에서 해당 지방자치단체의 조례로 정하는 금액을 부과한다.
  \item 허가권자는 이행강제금을 부과하기 전에 이행강제금을 부과ㆍ징수한다는 뜻을 미리 문서로써 계고하여야 한다.
  \item 허가권자는 위반 건축물에 대한 시정명령을 받은 자가 이를 이행하면 이미 부과된 이행강제금의 징수를 즉시 중지하여야 한다.
  \item 허가권자는 이행강제금 부과처분을 받은 자가 이행강제금을 납부기한까지 내지 아니하면 지방세 체납처분의 예에 따라 징수한다.
\end{choices}
\end{question}
\answer{4}
